\section{Parte 2: algoritmo exacto y voraz de aproximación}\label{sec:seleccion}

\subsection{Algoritmo exacto: árbol \emph{incluir / no incluir}}

Como \SEL{} es NP-difícil, el exacto se usa solo como referencia para medir al voraz en casos
pequeños (\texttt{seleccion\_exacta}).

\begin{algorithm}[H]
\caption{\textsc{SelecciónExacta}$(S_1,\dots,S_N,\ k)$ — \texttt{seleccion\_exacta}}\label{alg:exacto}
\SetKwFunction{Explorar}{Explorar}
\SetKwProg{Fn}{función}{:}{}
\Fn{\Explorar{$i$, $A$}}{
  \lIf(\tcp*[f]{caso 1: hoja}){$|A|=k$}{\Return{$\big(f(A),A\big)$}}
  \lIf(\tcp*[f]{caso 2: poda}){$N-i<k-|A|$}{\Return{$(-\infty,\emptyset)$}}
  $(v_1,A_1)\leftarrow$ \Explorar{$i+1$, $A\cup\{i\}$} \tcp*{caso 3a: incluir $i$}
  $(v_2,A_2)\leftarrow$ \Explorar{$i+1$, $A$} \tcp*{caso 3b: no incluir $i$}
  \Return{la mejor de las dos}\;
}
\Return{\Explorar{$0$, $\emptyset$}}\;
\end{algorithm}

\subsubsection*{Deducción 1: sin cupo, por recurrencia y sustitución}

Si se recorren todos los subconjuntos, cada llamada con $n$ nodos por decidir hace dos llamadas con
$n-1$ y trabajo constante $c$, y cada hoja evalúa $f$ desde cero con costo $\tau_f$:
\begin{align*}
T(n) &= 2T(n-1)+c,\qquad T(0)=\tau_f\\
     &= 2\big(2T(n-2)+c\big)+c = 2^2T(n-2)+c\,(2+1)\\
     &= 2^3T(n-3)+c\,(2^2+2+1)\\
     &\ \ \vdots\\
     &= 2^jT(n-j)+c\sum_{l=0}^{j-1}2^l = 2^jT(n-j)+c\,(2^j-1).
\end{align*}
Con $j=n$ se llega al caso base:
\begin{clave}
$T(N)=2^N\tau_f+c\,(2^N-1)=(\tau_f+c)\,2^N-c\ \in\ \Theta(2^N)$ si $\tau_f$ es constante; $\Theta(2^N\tau_f)$ en general.
\end{clave}

\subsubsection*{Deducción 2: con cupo $k$, recurrencia de Pascal}

Sea $L(n,r)$ el número de hojas (evaluaciones de $f$) cuando quedan $n$ nodos por decidir y faltan
$r$ por escoger. La rama \emph{incluir} gasta un lugar del cupo y la rama \emph{no incluir} no:
\[
  L(n,r)=L(n-1,r-1)+L(n-1,r),\qquad L(n,0)=1,\qquad L(n,r)=0\ \text{si } r>n .
\]
Es la regla de Pascal, $\binom{n}{r}=\binom{n-1}{r-1}+\binom{n-1}{r}$, con los mismos casos base; por
inducción sobre $n$, $L(n,r)=\binom{n}{r}$. El algoritmo evalúa $f$ \textbf{exactamente}
$\binom{N}{k}$ veces, y una prueba automática lo comprueba en 100 instancias. Como cada llamada interna lleva a una hoja y la profundidad es a lo sumo $N$:
\[
  T_{\text{exacto}}(N,k)\in\Omega\!\left(\tbinom{N}{k}\tau_f\right),\qquad
  T_{\text{exacto}}(N,k)\in O\!\left(\tbinom{N}{k}(N+\tau_f)\right),
\]
y el mejor y el peor caso coinciden, porque el recorrido no depende de los datos.

\begin{lema}
$\binom{N}{\floor{N/2}}\ge \dfrac{2^N}{N+1}$.
\end{lema}
\begin{proof}
Los $N+1$ números $\binom{N}{0},\dots,\binom{N}{N}$ suman $2^N$, y el mayor de ellos,
$\binom{N}{\floor{N/2}}$, es al menos su promedio.
\end{proof}

\noindent Con $k\approx N/2$ el exacto hace $\Omega(2^N/N)$ evaluaciones: es exponencial. Con $k$
fijo, $\binom{N}{k}\in\Theta(N^k)$, un polinomio cuyo grado crece con $k$; como $k$ es parte de la
entrada, el algoritmo no es polinomial en general. Por eso el programa solo lo ejecuta si
$\binom{N}{k}\le 500\,000$.

\subsection{Algoritmo voraz}

\begin{algorithm}[H]
\caption{\textsc{SelecciónVoraz}$(S_1,\dots,S_N,\ k)$ — \texttt{seleccion\_voraz}}\label{alg:voraz}
$A\leftarrow\emptyset$;\ $\mathit{cubierto}\leftarrow\emptyset$\;
\For{$\mathit{paso}\leftarrow 1$ \KwTo $k$}{
  \For{cada $j\notin A$}{
    $\Delta_j\leftarrow|S_j\setminus\mathit{cubierto}|$ \tcp*{una evaluación de la ganancia}
  }
  $j^*\leftarrow$ el $j$ de mayor $\Delta_j$ (en empate, el de mayor $p_j$; luego, el de menor índice)\;
  $A\leftarrow A\cup\{j^*\}$;\ $\mathit{cubierto}\leftarrow\mathit{cubierto}\cup S_{j^*}$\;
}
\Return{$A$}\;
\end{algorithm}

En el código, la ganancia es \texttt{(masks[j] \& \textasciitilde cubierto).bit\_count()}: las
categorías de $j$ que aún no están cubiertas, contadas como bits.

\paragraph{Desempate.} $p_j=u_j(1)-u_j(0)$ es la ganancia de la primera unidad de presupuesto del
nodo $j$; con $u_j(x)=w_j\ln(1+x)$, $p_j=w_j\ln2$, así que gana el de mayor coeficiente. El
desempate solo elige entre candidatos con el mismo $\Delta_j$ máximo, de modo que la sucesión de
valores $f_1,f_2,\dots$ y la garantía de abajo no cambian; comparar $p_j$ cuesta $O(1)$, así que el
conteo de evaluaciones tampoco cambia. Sí cambia \emph{qué} grupo se entrega cuando hay varios con
la misma cobertura: en el caso de la profesora, 33 grupos distintos cubren las 20 categorías, y su
utilidad tras el reparto va de 91.06 a 107.04. Desempatando por índice, el voraz entregaba uno de
92.58; desempatando por $p_j$, entrega uno de 99.58 (Sección~\ref{sec:resultados}).

\paragraph{Conteo de evaluaciones.} En el paso $i$ quedan $N-i+1$ candidatos:
\[
  T_{\text{voraz}}(N,k)=\sum_{i=1}^{k}(N-i+1)=kN-\frac{k(k-1)}{2}.
\]
Como $k\le N$, $k-1<N$ y $\frac{k(k-1)}{2}<\frac{kN}{2}$, así que $\frac{kN}{2}\le T_{\text{voraz}}\le kN$:
\begin{clave}
$T_{\text{voraz}}(N,k)\in\Theta(N k)$ evaluaciones y $\Theta(Nk\tau)$ de tiempo, con $\tau=O(|U|/w)$.
El mejor y el peor caso coinciden.
\end{clave}

\begin{teorema}[Nemhauser, Wolsey y Fisher, 1978]
Si $f$ es normalizada, monótona y submodular, el grupo $A_g$ del voraz cumple
$f(A_g)\ge\left(1-\frac1e\right)f(A^*)\approx 0.632\,f(A^*)$.
\end{teorema}
\begin{proof}[Idea de la demostración]
Sea $\mathrm{OPT}=f(A^*)$ y $f_i$ el valor tras $i$ pasos. Por monotonía,
$\mathrm{OPT}\le f(A_i\cup A^*)$, con $A_i$ el grupo del voraz tras $i$ pasos. Por submodularidad, esa diferencia se reparte en a lo sumo $k$
aportes de nodos de $A^*$, así que alguno aporta al menos $(\mathrm{OPT}-f_i)/k$, y el voraz toma
uno al menos igual de bueno: $\mathrm{OPT}-f_{i+1}\le(1-\frac1k)(\mathrm{OPT}-f_i)$. Tras $k$ pasos,
$\mathrm{OPT}-f_k\le(1-\frac1k)^k\mathrm{OPT}\le\mathrm{OPT}/e$, porque $1-y\le e^{-y}$.
\end{proof}

\paragraph{Aplicación.} Las hipótesis son las del Lema~\ref{lem:valida}. En cada caso pequeño el
programa calcula $f(A_g)/f(A^*)$ y verifica que sea $\ge 1-1/e$; las pruebas lo comprueban en 300
instancias aleatorias. La cota no se puede mejorar en tiempo polinomial salvo que P$=$NP
(Feige, 1998), así que el voraz alcanza la mejor garantía posible.

\paragraph{El voraz no siempre es óptimo.} Con $S_1=\{1,2,3,4\}$, $S_2=\{1,2,5\}$, $S_3=\{3,4,6\}$
y $k=2$ (caso \texttt{casos/ejemplos/ej00\_tres\_nodos.json}), el voraz toma $S_1$ y luego $S_2$
o $S_3$ (los dos aportan un elemento nuevo), con $f=5$, pero $\{S_2,S_3\}$ cubre 6. La razón $5/6\approx0.833$ cumple la garantía.

\subsection{Resumen de funciones de eficiencia}

\begin{center}
\begin{tabularx}{\linewidth}{lXc}
\cabecera \cab{Algoritmo} & \cab{Función de eficiencia} & \cab{¿Mejor $=$ peor?} \\
Reparto voraz (P1) & peor $\Theta(N+B\log N)$; mejor $\Theta(N+B)$ & No \\
Selección exacta con cupo (P2) & $\binom{N}{k}$ evaluaciones; $O\big(\binom{N}{k}(N+\tau_f)\big)$ & Sí \\
Selección exacta sin cupo (P2) & $(\tau_f+c)2^N-c\in\Theta(2^N\tau_f)$ & Sí \\
Selección voraz (P2) & $kN-\frac{k(k-1)}{2}=\Theta(Nk)$ evaluaciones; $\Theta(Nk\tau)$ tiempo & Sí \\
Solución base completa & $\Theta(Nk\tau)+O(k+B\log k)$, más leer la entrada en $O(N+\sum_j|S_j|)$ & No \\
\bottomrule
\end{tabularx}
\end{center}
